\begin{pastexam}[label={px:gk3-113-ds-04}]{高考三級資訊處理 113 年 第四題}
（申論；25 分）

堆疊（Stack）與佇列（Queue）是常見的資料結構，請回答下列問題：
\begin{itemize}[leftmargin=2.2em]
\item[㈠] 利用雙向佇列（Deque）循序輸入 1, 2, 3, 4, 5, 6, 7，請問能否得到 5174236 的輸出排列？並說明其過程
或理由。（10 分）
\item[㈡] 若有 1, 2, 3, 4 四個數字要依序 Push 進堆疊，再於任意時間點 Pop 出堆疊，請列出可能的輸出組合。（15 分）
\end{itemize}
\end{pastexam}

\begin{pastsoln}
答案：㈠ 不能。 ㈡ 共 $14$ 種：1234、1243、1324、1342、1432、2134、2143、2314、2341、2431、3214、
3241、3421、4321。

\textbf{㈠} 雙向佇列兩端都能放入、都能取出（資結書第~6.2~節）。 $1 \sim 7$ 依序放入，任何時候可以從
任一端取出一個輸出。 兩端都不受限制時做不到，限制輸入端或輸出端的版本能做的更少，當然也做不到。

\textbf{理由。} 先看一個性質：只從兩端放入時，以最早放進去、還在裡面的那個數為中心，往左、往右
兩邊都是越外面的數越晚放入，也就是越大。

\begin{enumerate}
\item 第一個輸出是 $5$，所以輸出之前 $1, 2, 3, 4, 5$ 都已經放進去了。 $5$ 是剛放入、在某一端的數，取出它
沒有問題。
\item 第二個輸出是 $1$。 這時裡面有 $1, 2, 3, 4$（也可能先放了 $6$、$7$，但它們都在最外面，不影響下面的
論證）。 $1$ 最早放入，位於中心；要從端點取出 $1$，$1$ 的某一邊必須是空的，也就是 $2, 3, 4$ 全部在
$1$ 的同一邊。 由上面的性質，那一邊由內往外是 $2, 3, 4$，所以取出 $1$ 之後，$2, 3, 4$ 是相鄰的一段，
順序為 $2, 3, 4$（$2$ 在一端那一側的內部，$4$ 在外側）。
\item 接下來要輸出 $7, 4, 2, 3, 6$。 $6, 7$ 只能放在兩端的最外面。 輸出 $4$ 時，$4$ 要在端點，所以 $6$ 不能
在 $4$ 的外側，只能在 $2$ 那一側：排列是 $6, 2, 3, 4$（或左右相反）。 取出 $4$ 後剩 $6, 2, 3$，兩端是
$6$ 和 $3$，$2$ 夾在中間，下一個輸出不可能是 $2$。
\end{enumerate}
所以 $5174236$ 不能由雙向佇列得到。

\textbf{㈡} $n$ 個數依序 push、任意時間 pop，能得到的輸出排列共有卡塔蘭數 $\frac{1}{n+1}\binom{2n}{n}$ 種
（資結書定理~6.2），$n = 4$ 時是 $\frac{1}{5}\binom{8}{4} = 14$ 種。 依第一個輸出的數分類列舉：
\begin{itemize}
\item \textbf{先輸出 1}：push 1、pop 1，剩下 $2, 3, 4$ 自己做，有 $5$ 種：1234、1243、1324、1342、1432。
\item \textbf{先輸出 2}：push 1、2，pop 2，此時 $1$ 在堆疊底。 $1$ 可以接著出（2134、2143），也可以等 3 出來
之後出（2314），或等 4 出來之後出（2341、2431）。 共 $5$ 種。
\item \textbf{先輸出 3}：push 1、2、3，pop 3，堆疊裡是 $1, 2$（$2$ 在上），$2$ 必須比 $1$ 先出。 可能的有
3214、3241、3421，共 $3$ 種。 3124、3142、3412 都不行（$1$ 在 $2$ 之前出）。
\item \textbf{先輸出 4}：四個都 push 進去才 pop，只能 4321，$1$ 種。
\end{itemize}
合計 $5 + 5 + 3 + 1 = 14$ 種。 判斷某個排列行不行，可用資結書定理~6.1：不能出現 $p_j < p_k < p_i$
（$i < j < k$）的「3-1-2」型態。 例如 3142 裡的 $3, 1, 2$ 就是，所以不行；全部 $24$ 種排列扣掉含這種型態
的 $10$ 種，正好剩 $14$ 種。
\end{pastsoln}
